\documentclass[11pt]{amsart}
\usepackage[T1]{fontenc}
\usepackage{lmodern}
\usepackage{amsmath,amssymb,amsthm}
\usepackage[margin=1.1in]{geometry}
\usepackage{booktabs}
\usepackage{hyperref}

\newtheorem{theorem}{Theorem}
\newtheorem{lemma}[theorem]{Lemma}
\newtheorem{proposition}[theorem]{Proposition}
\newtheorem{corollary}[theorem]{Corollary}
\newtheorem{conjecture}[theorem]{Conjecture}
\newtheorem{question}[theorem]{Question}
\theoremstyle{definition}
\newtheorem{definition}[theorem]{Definition}
\newtheorem{remark}[theorem]{Remark}
\newtheorem{example}[theorem]{Example}

\newcommand{\Ev}{\mathrm{Ev}}
\newcommand{\cA}{\mathcal A}
\newcommand{\cD}{\mathcal D}

\title[Multi-pair separation]{Separating many pairs in Lov\'asz decompositions:\\
classes, the leaf merge, and the Case~I lemma}
\author{Justin Williams}
\date{November 2026 (working draft)}

\begin{document}
\begin{abstract}
A \emph{Lov\'asz decomposition} of a graph with at most one vertex of positive even degree is a
partition of its edges into paths in which every vertex $v$ is the end of exactly $d(v)\bmod 2$
paths; Lov\'asz proved that one exists. We study the \emph{multi-pair separation conjecture} (MSEP):
if all degrees are odd and $\{c_1,a_1\},\dots,\{c_t,a_t\}$ are disjoint pairs at distance at least three,
then some Lov\'asz decomposition has no path containing both vertices of a pair. MSEP would give
Gallai's conjecture, with every orientation of the matching realisable, for every graph whose even
vertices induce a subgraph with a perfect matching none of whose edges lies in a triangle; this class
contains every triangle-free Eulerian graph with a perfect matching. We re-check this reduction.

We do not prove MSEP. We prove the following.
\begin{itemize}
\item \emph{Prescribed end edges.} If all degrees are odd, then for every vertex $x$ and every edge
$xu$ there is a Lov\'asz decomposition in which $xu$ is the terminal edge of the path ending at $x$;
moreover every path can be made to meet a given packed set containing $x$ at most once.
\item \emph{Classes.} MSEP is the case of pairs of a \emph{class separation} conjecture (CSEP) for
partitions of the vertex set into packed sets, equivalently for proper colourings of the square of
the graph. Lov\'asz's Case~II costs nothing for CSEP: the new leaf joins the class of the split
vertex, and the lift is a merge (\emph{leaf merge}). So CSEP reduces to graphs with exactly one
vertex of positive even degree.
\item \emph{The Case~I lemma.} For such graphs, Lov\'asz's Case~I with a trivial Fan chain is
equivalent to prescribing, at a neighbour $u$ of the even vertex $z$, the terminal edge $uz$. We
isolate this single statement, prove a certified version of the step, and a rotation lemma that
removes a cycle without creating an even vertex.
\item \emph{Sparse classes.} If every vertex has at most one neighbour in the union of the classes,
CSEP holds, even for overlapping classes. This contains the separation theorem for packed sets.
\item \emph{Covers.} CSEP holds for the fibres of a covering map onto a graph with all degrees odd.
\end{itemize}
The proofs are computation-free. The Lean library formalises the reduction to the Gallai class, the leaf
merge, the prescribed end edge, the Case~II reduction, the certified Case~I step and the sparse-class
theorem. Separately labelled
computations give the evidence: no counterexample to MSEP among all $36{,}024{,}731$ distance-three
perfect pairings of the $4{,}060$ connected cubic graphs of order $16$, $8.26$ million sampled
pairings at order $18$, and random graphs to order $20$; and no counterexample to CSEP for minimum
colourings of the square. We also show that every natural strengthening we tried of the inductive
invariant is false, and that certified local moves alone cannot prove MSEP.
\end{abstract}
\maketitle

\section{Introduction}\label{sec:intro}

All graphs are finite and simple. A \emph{path decomposition} of a graph $K$ is a partition of $E(K)$
into paths, each with at least one edge. For such a decomposition $\cD$ and a vertex $v$, $b_\cD(v)$ is
the number of paths of $\cD$ with an end at $v$; always $b_\cD(v)\equiv d_K(v)\pmod 2$. A vertex has
\emph{positive even degree} if $d(v)>0$ is even. Following~\cite{RD} we call $\cD$ a \emph{Lov\'asz
decomposition} if $b_\cD(v)=d_K(v)\bmod 2$ for every $v$. Lov\'asz~\cite{L68} proved:

\begin{theorem}[Lov\'asz]\label{thm:lovasz}
If $K$ has at most one vertex of positive even degree, it has a Lov\'asz decomposition.
\end{theorem}

A set $X$ of vertices is \emph{packed} if its vertices are pairwise at distance at least three: no two
are adjacent and no two have a common neighbour. Equivalently, the closed neighbourhoods $N[x]$,
$x\in X$, are pairwise disjoint. The route-descent manuscript~\cite{RD} proves the \emph{separation
theorem}: if $K$ has at most one vertex of positive even degree and $X$ is packed, then some Lov\'asz
decomposition has every path meeting $X$ in at most one vertex. The case $|X|=2$ separates two vertices.

\begin{conjecture}[MSEP, multi-pair separation]\label{conj:msep}
Let $K$ have all degrees odd, and let $\{c_1,a_1\},\dots,\{c_t,a_t\}$ be pairwise disjoint pairs with
$\mathrm{dist}_K(c_i,a_i)\ge3$. Then $K$ has a Lov\'asz decomposition in which no path contains both
$c_i$ and $a_i$, for any $i$.
\end{conjecture}

Conjecture~\ref{conj:msep} is stated in~\cite{RD}, with evidence. Its interest is
Proposition~\ref{prop:reduction}: it would prove Gallai's conjecture for a class of graphs in which the
even vertices may induce a dense subgraph.

\subsection*{Results}
\begin{enumerate}
\item Proposition~\ref{prop:reduction} re-checks the reduction of a Gallai class to MSEP.
\item Section~\ref{sec:classes} introduces \emph{class separation} (Conjecture~\ref{conj:csep}). The
vertex set is partitioned into packed sets, and every path should meet every class at most once. MSEP is
the case of classes of size at most two. Partitions into packed sets are the fibres of locally
injective homomorphisms, or the colour classes of proper colourings of $K^2$. We prove the conjecture
for fibres of coverings (Proposition~\ref{prop:cover}).
\item Section~\ref{sec:merge}: the \emph{leaf merge} (Lemma~\ref{lem:merge}) and two consequences.
Theorem~\ref{thm:endedge}: in a graph with all degrees odd, any single edge can be prescribed as the
terminal edge at its end vertex, together with separation of a packed set. Theorem~\ref{thm:caseII}:
Lov\'asz's Case~II is free for class separation, so CSEP reduces to graphs with exactly one vertex of
positive even degree.
\item Section~\ref{sec:caseI}: the remaining case. The step of Lov\'asz's induction at the even vertex
$z$ succeeds with a trivial Fan chain exactly when some neighbour $u$ of $z$ can be given the terminal
edge $uz$ (Proposition~\ref{prop:EE}). We prove a certified form of the step (Lemma~\ref{lem:caseI}) and
a rotation lemma for cycles (Lemma~\ref{lem:rotation}).
\item Section~\ref{sec:sparse}, Theorem~\ref{thm:sparse}: class separation holds, for arbitrary
overlapping families of packed sets, when every vertex has at most one neighbour in their union. This
contains the separation theorem.
\item Section~\ref{sec:limits} (computational, labelled): the evidence for MSEP and CSEP, and the
failure of the natural strengthenings of the inductive invariant.
\item Section~\ref{sec:lean}: the Lean formalisation.
\end{enumerate}

\section{Preliminaries}\label{sec:prelim}

For a path decomposition $\cD$ and a vertex $v$ with $b_\cD(v)=1$, $Q_v$ is the path of $\cD$ ending at
$v$. If $v$ is an end of a path $P$, the \emph{terminal edge of $P$ at $v$} is the edge of $P$ at $v$.
We use Fan's lemma in the form proved in~\cite{RD}.

\begin{lemma}[Fan's lemma with vertex sets~\cite{Fan05,RD}]\label{lem:fan}
Let $\cD$ be a path decomposition of $K$, let $t\ne u$ be non-adjacent, and write
$Z(u)=\{w\in N(u): b_\cD(w)=0\}$. If $b_\cD(t)>|Z(u)|$, then $K+tu$ has a path decomposition $\cD^+$
with $b_{\cD^+}=b_\cD-\mathbf 1_t+\mathbf 1_u$ and a bijection $\varphi:\cD\to\cD^+$ such that
$V(\varphi(Q))=V(Q)$ for every $Q$ except one, $Q_0$, with $V(\varphi(Q_0))=V(Q_0)\cup\{u\}$.
\end{lemma}

The path $Q_0$ of Lemma~\ref{lem:fan} is the path to which the last edge of the Fan chain is appended.
It ends at a neighbour of $u$ or at $t$. In the \emph{trivial chain} case, $u\notin V(Q_t)$ for a path
$Q_t$ ending at $t$; then $\cD^+$ is obtained by appending $tu$ to $Q_t$.

\begin{theorem}[separation theorem for packed sets~\cite{RD}]\label{thm:sep}
If $K$ has at most one vertex of positive even degree and $X$ is packed, then $K$ has a Lov\'asz
decomposition in which every path meets $X$ in at most one vertex.
\end{theorem}

\emph{The leaf graph.} For an edge $xu$ of $K$ let $H=H(K;x,u)$ be $K-xu$ together with a new vertex
$\ell$ and the edge $u\ell$. Let $f:V(H)\to V(K)$ be the identity on $V(K)$ and send $\ell$ to $x$. In
$H$ the vertex $x$ has degree $d_K(x)-1$, the leaf $\ell$ has degree one, and all other degrees are
those of $K$. This is the graph of Lov\'asz's Case~II.

\section{The reduction to a Gallai class}\label{sec:reduction}

For a graph $G$ let $\Ev$ be its set of vertices of positive even degree. If $N$ is a perfect matching
of $G[\Ev]$, an \emph{orientation} of $N$ names one end $a_i$ (the head) and the other $c_i$ (the tail)
of each edge $c_ia_i$ of $N$. A path decomposition of $G$ \emph{realises} the orientation if every head
is an end of exactly two paths, every tail of none, and every vertex of odd degree of exactly one.

\begin{proposition}\label{prop:reduction}
Let $G$ be a graph without isolated vertices of order $n$, let $N$ be a perfect matching of $G[\Ev]$,
and let $F=G-N$.
\begin{enumerate}
\item[(a)] Every vertex has odd degree in $F$.
\item[(b)] For $ca\in N$, $\mathrm{dist}_F(c,a)\ge3$ if and only if $ca$ lies in no triangle of $G$.
\item[(c)] If $F$ has a Lov\'asz decomposition in which no path contains both ends of an edge of $N$, then
every orientation of $N$ is realised by a path decomposition of $G$ into $n/2$ paths. In particular
$G$ satisfies Gallai's bound $\lceil n/2\rceil$.
\end{enumerate}
So Conjecture~\ref{conj:msep} implies Gallai's conjecture, with every orientation of $N$ realisable, for
every graph whose even vertices induce a subgraph having a perfect matching none of whose edges lies in
a triangle.
\end{proposition}

\begin{proof}
(a) Every vertex of $\Ev$ loses exactly one edge, and the other vertices keep their odd degree.

(b) As $G$ is simple, $ca\notin E(F)$, so $\mathrm{dist}_F(c,a)\ge2$, with equality iff some $w$ is
adjacent in $F$ to $c$ and $a$; then $cwa$ is a triangle of $G$. Conversely, if $caw$ is a triangle of
$G$, the edges $cw$ and $aw$ are not in the matching $N$, which already covers $c$ and $a$ by $ca$; so
they lie in $F$.

(c) Let $\cD$ be the decomposition of $F$. By (a) every vertex is the end of exactly one path of $\cD$,
so $|\cD|=n/2$, and $n$ is even. Fix an orientation. For each path $Q\in\cD$ and each end $s$ of $Q$
that is a tail $c_i$, attach the edge $c_ia_i$ at $s$. Since $Q$ contains $c_i$, it does not contain
$a_i$. If both ends of $Q$ are tails $c_i$ and $c_j$, then $i\ne j$, the heads $a_i\ne a_j$ are both
off $Q$, and $a_ic_iQc_ja_j$ is a path. So every extended $Q$ is a path. Every edge of $N$ is attached
exactly once, at the unique path ending at its tail. A tail loses its only end. A head keeps its end in
$\cD$ and gains the end of the attached edge. A vertex of odd degree of $G$ keeps its single end. The
number of paths is $|\cD|=n/2$.
\end{proof}

\begin{remark}[on novelty]\label{rem:novelty}
Gallai's conjecture is open for Eulerian graphs; it is known for maximum degree at most four~\cite{FK}
and five~\cite{BP19}, and~\cite{CW25} proves $3n/5$ for Eulerian graphs whose triangles are pairwise at
distance at least three. Pyber~\cite{Pyb96} (even subgraph a forest), Fan~\cite{Fan05} (blocks of the
even subgraph triangle-free with maximum degree at most three) and Botler and Sambinelli~\cite{BS21}
impose sparsity on $G[\Ev]$. The class of Proposition~\ref{prop:reduction} imposes only a perfect
matching avoiding triangles. It contains every triangle-free Eulerian graph with a perfect matching, and
graphs whose even subgraph has large degrees and many cycles. We found no source treating it.
\end{remark}

\section{Class separation}\label{sec:classes}

\begin{definition}\label{def:classes}
A \emph{class family} of $K$ is a family $\cA$ of pairwise disjoint packed sets (\emph{classes}). A path
decomposition is \emph{rainbow} for $\cA$ if every path meets every class in at most one vertex. For a
vertex $v$, $\cA(v)$ denotes the class containing $v$, if any.
\end{definition}

\begin{conjecture}[CSEP, class separation]\label{conj:csep}
If $K$ has at most one vertex of positive even degree and $\cA$ is a class family of $K$, then $K$ has
a Lov\'asz decomposition that is rainbow for $\cA$.
\end{conjecture}

For classes of size at most two and all degrees odd this is Conjecture~\ref{conj:msep}. For one class it is
Theorem~\ref{thm:sep}. Adding singleton classes changes nothing, so we may take $\cA$ to be a partition of
$V(K)$.

\emph{Reformulations.} A partition of $V(K)$ into packed sets is the same as the partition into fibres of
a map $\varphi:V(K)\to V(L)$ that is a homomorphism onto a multigraph $L$ (the quotient) and is injective
on every closed neighbourhood. These are the \emph{locally injective homomorphisms}, or partial covers.
Equivalently, the classes are the colour classes of a proper colouring of the square $K^2$. So
Conjecture~\ref{conj:csep} says: \emph{for every proper colouring of $K^2$, some Lov\'asz decomposition
of $K$ has all its paths rainbow.} Paths are automatically rainbow on vertices within distance two of each
other, and the conjecture asks that the colouring never forces two far-apart vertices of the same colour
onto one path.

\begin{remark}\label{rem:overlap}
For overlapping families of packed sets the statement is false. The union of three perfect pairings at
distance three of a cubic graph of order $12$ can admit no rainbow Lov\'asz decomposition
(Section~\ref{sec:limits}). Disjointness matters.
\end{remark}

\begin{proposition}[covers]\label{prop:cover}
Let $\varphi:K\to L$ be a covering map, that is, a surjective homomorphism that maps $N(v)$ bijectively
onto $N(\varphi(v))$ for every $v$. If every vertex of $L$ has odd degree, then the fibres of $\varphi$ form
a class family of $K$, and $K$ has a Lov\'asz decomposition that is rainbow for it.
\end{proposition}

\begin{proof}
Two vertices of a fibre are not adjacent, since $L$ has no loops. They have no common neighbour $w$,
since $\varphi$ is injective on $N(w)$. So fibres are packed. Let $\cD_L$ be a Lov\'asz decomposition of
$L$ (Theorem~\ref{thm:lovasz}). For a path $P=p_0p_1\cdots p_k$ of $\cD_L$ and $v\in\varphi^{-1}(p_0)$
there is a unique walk $v=v_0v_1\cdots v_k$ with $\varphi(v_i)=p_i$: $v_{i+1}$ is the unique neighbour of
$v_i$ in $\varphi^{-1}(p_{i+1})$. Its vertices are distinct, since their images are, so it is a path, and it
meets every fibre at most once. Every edge $vw$ of $K$ lies on exactly one such lift: $\varphi(v)\varphi(w)$
lies on exactly one path $P$ of $\cD_L$, and the lift of $P$ through $v$ in the direction of $w$ is
determined. The lifts of $P$ give a bijection between the fibres over its two ends. Each vertex $v$ of $K$
is therefore the end of exactly one lift, of the path of $\cD_L$ ending at $\varphi(v)$. As
$d_K(v)=d_L(\varphi(v))$ is odd, the lifts form a Lov\'asz decomposition of $K$.
\end{proof}

For a $2$-fold covering $K\to L$ of a graph with all degrees odd, Propositions~\ref{prop:cover}
and~\ref{prop:reduction} give Gallai's bound, with every orientation realisable, for $K$ plus the
matching of fibres. This easy class shows that Conjecture~\ref{conj:csep} interpolates between covers,
where it is a lifting statement, and Theorem~\ref{thm:sep}.

\section{The leaf merge}\label{sec:merge}

\begin{lemma}[leaf merge]\label{lem:merge}
Let $xu$ be an edge of $K$, let $H=H(K;x,u)$ and $f$ be as in Section~\ref{sec:prelim}, and let
$\cD_H$ be a path decomposition of $H$ in which no path contains both $x$ and $\ell$. Then
$\cD=\{f(P):P\in\cD_H\}$ is a path decomposition of $K$, $V(f(P))=f(V(P))$ for every $P$,
$b_\cD(x)=b_{\cD_H}(x)+b_{\cD_H}(\ell)$, and $b_\cD(w)=b_{\cD_H}(w)$ for $w\ne x$.
\end{lemma}

\begin{proof}
The map $f$ sends the edge $u\ell$ to $ux$ and fixes the other edges of $H$, which are the edges of
$K-xu$. So $f$ is a homomorphism $H\to K$ that is a bijection on edges. It identifies only $x$ and $\ell$,
so it is injective on $V(P)$ for every $P\in\cD_H$, and $f(P)$ is a path with the same number of edges,
whose ends are the images of the ends of $P$. The images partition $E(K)$, and the end counts are as
stated.
\end{proof}

\begin{lemma}\label{lem:packedlift}
Let $A$ be packed in $K$. Then $A$ is packed in $H$ if $x\notin A$, and $A\cup\{\ell\}$ is packed in $H$
if $x\in A$.
\end{lemma}

\begin{proof}
$H-\ell$ is a subgraph of $K$, so two vertices of $A$ remain at distance at least three. Let $x\in A$
and $a\in A$. The only neighbour of $\ell$ is $u$, and $u\notin A$ since $u$ is adjacent to $x\in A$. So
$\ell\not\sim a$. A common neighbour of $\ell$ and $a$ would be $u$, with $ua\in E(H)$. Then $a\ne x$, as
$xu\notin E(H)$, and $u$ would be a common neighbour of $x$ and $a$ in $K$, which is impossible.
\end{proof}

\begin{theorem}[prescribed end edge]\label{thm:endedge}
Let $K$ have no vertex of positive even degree, let $X$ be packed, let $x\in X$, and let $xu\in E(K)$.
Then $K$ has a Lov\'asz decomposition in which every path meets $X$ at most once and in which the path
ending at $x$ has terminal edge $xu$. In particular (take $X=\{x\}$), every edge $xu$ is the terminal
edge at $x$ of some Lov\'asz decomposition.
\end{theorem}

\begin{proof}
In $H=H(K;x,u)$ only $x$ can have positive even degree, namely $d_K(x)-1$. By
Lemma~\ref{lem:packedlift}, $X\cup\{\ell\}$ is packed in $H$. Theorem~\ref{thm:sep} gives a Lov\'asz
decomposition $\cD_H$ of $H$ in which every path meets $X\cup\{\ell\}$ at most once; in particular no
path contains $x$ and $\ell$. Let $\cD$ be as in Lemma~\ref{lem:merge}. Then
$b_\cD(x)=(d_K(x)-1\bmod2)+1=1$, as $d_K(x)$ is odd, and $b_\cD(w)=d_K(w)\bmod2$ for $w\ne x$; so $\cD$ is a
Lov\'asz decomposition. A path $f(P)$ meets $X$ in $f(V(P)\cap(X\cup\{\ell\}))$, so at most once. The
path $P_\ell$ of $\cD_H$ ending at the leaf contains its only edge $u\ell$. So $f(P_\ell)$ ends at $x$ and
contains $ux$, which is its terminal edge at $x$ because $x$ occurs on it once.
\end{proof}

Theorem~\ref{thm:endedge} is the counterpart of the leaf trick by which~\cite{RD} proves $L_1$. There a
leaf is attached at $c$ and separated from $a$. Here an edge is detached into a leaf, which is separated
from its old end and merged back. We did not find Theorem~\ref{thm:endedge} in the literature.
Prescribing end edges at several vertices of a packed set at once is open
(Question~\ref{q:pe}).

\begin{theorem}[Case~II is free]\label{thm:caseII}
Let $K$ have no vertex of positive even degree, let $\cA$ be a class family of $K$, and let $xu\in E(K)$.
In $H=H(K;x,u)$ let $\cA_H$ consist of the classes $A\in\cA$ with $x\notin A$, together with
$\cA(x)\cup\{\ell\}$ if $x$ lies in a class, and $\{x,\ell\}$ otherwise. Then $\cA_H$ is a class family
of $H$, and if $H$ has a Lov\'asz decomposition rainbow for $\cA_H$, then $K$ has one rainbow for $\cA$.
Consequently, Conjecture~\ref{conj:csep} for all graphs with exactly one vertex of positive even degree
implies it for all graphs with none.
\end{theorem}

\begin{proof}
The sets of $\cA_H$ are packed by Lemma~\ref{lem:packedlift} and pairwise disjoint, since $\ell$ lies in
exactly one of them. Let $\cD_H$ be rainbow for $\cA_H$. Some class of $\cA_H$ contains $x$ and $\ell$, so
no path contains both, and Lemma~\ref{lem:merge} applies. The end counts are those of the proof of
Theorem~\ref{thm:endedge}. A path $f(P)$ meets a class $A\in\cA$ in $f(V(P)\cap A')$, where $A'$ is the
class of $\cA_H$ corresponding to $A$; so at most once. For the consequence: if some vertex $x$ has degree
at least three, $H(K;x,u)$ has exactly one vertex of positive even degree, namely $x$. Otherwise every
degree is at most one, every path of every decomposition is a single edge, and a single edge meets a
packed set at most once.
\end{proof}

\section{The Case~I lemma}\label{sec:caseI}

By Theorem~\ref{thm:caseII}, Conjecture~\ref{conj:csep} comes down to graphs $K$ with exactly one vertex
$z$ of positive even degree. Lov\'asz's induction deletes an edge $uz$, decomposes $K'=K-uz$ (whose only
positive even vertex is $u$), and restores $uz$ by Lemma~\ref{lem:fan} with target $u$. The path that
gains $u$ ends at $z$ or at a neighbour of $u$, and it may contain another vertex of the class of $u$.
When $u$ is in no class this is harmless, and this is the proof of Theorem~\ref{thm:sep}. When all
vertices are in non-trivial classes, as for MSEP with a perfect pairing, every target is in a class.

\begin{proposition}[Case~I is an end-edge prescription]\label{prop:EE}
Let $K$ have exactly one vertex $z$ of positive even degree, let $u\in N(z)$, $K'=K-uz$, and let $\cA$ be
a class family of $K$. The following are equivalent.
\begin{enumerate}
\item[(i)] $K$ has a Lov\'asz decomposition, rainbow for $\cA$, in which $uz$ is the terminal edge at $u$
of the path ending at $u$.
\item[(ii)] $K'$ has a Lov\'asz decomposition, rainbow for $\cA$, in which the path ending at $z$ avoids
$u$ and every vertex of $\cA(u)$.
\end{enumerate}
\end{proposition}

\begin{proof}
In $K'$ the vertex $z$ is odd and $u$ is even or isolated. Given (ii), append $uz$ to the path $Q_z$ of
$K'$ ending at $z$. This is a path because $u\notin Q_z$; $z$ loses its end and $u$ gains one, so the
result is a Lov\'asz decomposition of $K$. It meets $\cA(u)$ only in $u$ and every other class as $Q_z$
did. Given (i), let $P=Rzu$ be the path ending at $u$ with terminal edge $uz$; $z$ is interior, as it
has no end. Replacing $P$ by $Rz$ gives a Lov\'asz decomposition of $K'$: $u$ loses its only end and $z$
gains one. $Rz$ is the path ending at $z$; it avoids $u$, and it avoids $\cA(u)$ because $P$ is rainbow.
\end{proof}

So Lov\'asz's induction proves Conjecture~\ref{conj:csep} as soon as, at every graph with one positive even
vertex $z$, \emph{some} neighbour $u$ satisfies (i). We state the form we tested.

\begin{conjecture}[EE, end edge at the even vertex]\label{conj:EE}
If $K$ has exactly one vertex $z$ of positive even degree and $\cA$ is a class family, then for every
$u\in N(z)$, $K$ has a Lov\'asz decomposition, rainbow for $\cA$, in which $uz$ is the terminal edge at $u$.
\end{conjecture}

Conjecture~\ref{conj:EE} implies Conjecture~\ref{conj:csep} directly for one positive even vertex, and then
for none by Theorem~\ref{thm:caseII}. Compare Theorem~\ref{thm:endedge}, which is the analogue without an
even vertex and is a theorem. Condition (ii) is an $L_1$-type statement: the end path of $z$ must avoid a
set at distance at least two from $z$ in $K'$ (as $z\sim u$ in $K$ and the other vertices of $\cA(u)$ are
at distance at least three from $u$). The leaf trick that proves $L_1$ in~\cite{RD}
attaches a leaf at $z$, which here would create a second even vertex; and $L_1$ is false in graphs with
an extra even vertex adjacent to the avoided vertex (project experiment E20). In (ii) the avoided set contains the
even vertex $u$ itself, and we found no failure (Section~\ref{sec:limits}).

The next lemma certifies the step when distances allow it.

\begin{lemma}[certified Case~I]\label{lem:caseI}
Let $K$, $z$, $u$, $K'$ be as in Proposition~\ref{prop:EE}, and let $\cA$ be a family of packed sets of $K$
(not necessarily disjoint). Suppose $K'$ has a Lov\'asz decomposition $\cD'$ meeting $\{u,z\}$ at most once
per path, meeting $A\cup\{z\}$ at most once per path for every $A\in\cA$ containing $u$, and meeting every
other $A\in\cA$ at most once per path. Then appending $uz$ to the path of $\cD'$ ending at $z$ gives a
Lov\'asz decomposition of $K$ meeting every $A\in\cA$ at most once per path.
\end{lemma}

\begin{proof}
Let $Q$ be the path of $\cD'$ ending at $z$; $u\notin Q$ because $Q$ meets $\{u,z\}$ only in $z$. The end
counts are as in Proposition~\ref{prop:EE}. The paths other than $Q+zu$ are unchanged. Let $A\in\cA$. If
$u\notin A$, the path $Q+zu$ meets $A$ as $Q$ does. If $u\in A$, then $z\notin A$ ($z\sim u$), and $Q$ meets
$A\cup\{z\}$ only in $z$, so $Q\cap A=\emptyset$ and $Q+zu$ meets $A$ only in $u$.
\end{proof}

The hypothesis of Lemma~\ref{lem:caseI} is a rainbow condition for the family in which $z$ has joined the
class of $u$. It is meaningful only if $\cA(u)\cup\{z\}$ and $\{u,z\}$ are packed in $K'$, that is, if $z$
is at distance at least three in $K'$ from $u$ and from every vertex of $\cA(u)$. The family is no longer
disjoint, and the induction that uses Lemma~\ref{lem:caseI} at every step must carry overlapping families;
these grow along the trails of Lov\'asz's induction, and we show in Section~\ref{sec:limits} that the
distance conditions eventually fail. The same holds for the following move, which never creates an even
vertex.

\begin{lemma}[rotation]\label{lem:rotation}
Let $K$ have all degrees odd, let $C=c_0c_1\cdots c_{r-1}c_0$ be a cycle, and let $K_C=K-E(C)$ (all degrees
odd). Indices are modulo $r$. Let $\cD_C$ be a Lov\'asz decomposition of $K_C$, rainbow for a class family
$\cA$, such that for every $i$ the path $Q_i$ ending at $c_i$ avoids $c_{i+1}$ and $\cA(c_{i+1})$, and such
that $\cA(c_{i+1})\ne\cA(c_{j+1})$ or both are undefined whenever $Q_i=Q_j$ with $i\ne j$. Then extending
every $Q_i$ by the edge $c_ic_{i+1}$ at $c_i$ gives a Lov\'asz decomposition of $K$, rainbow for $\cA$.
\end{lemma}

\begin{proof}
Each $c_i$ is an end of exactly one path $Q_i$ of $\cD_C$. Extending $Q_i$ at $c_i$ by $c_ic_{i+1}$ gives a
path, since $c_{i+1}\notin Q_i$. If $Q_i=Q_j$ for some $j\ne i$, the path is extended at both ends by
$c_{i+1}\ne c_{j+1}$, both off $Q_i$, so it is still a path. Every edge of $C$ is used once. The vertex $c_i$
loses its end and gains the end of the extension of $Q_{i-1}$, so it still has one; other vertices are
unchanged. The extension of $Q_i$ meets $\cA(c_{i+1})$ only in $c_{i+1}$, and when both ends are extended
the two added vertices lie in different classes, so every extended path is rainbow.
\end{proof}

A leaf-to-leaf path $p_0p_1\cdots p_m$ admits the same treatment, with $Q_{p_i}$ extended for $1\le i\le
m-1$ and $p_0p_1$ added as a new path. The conditions of Lemma~\ref{lem:rotation} can be certified, as in
Lemma~\ref{lem:caseI}, by letting $c_i$ join the class of $c_{i+1}$ in $K_C$.

\section{Sparse classes}\label{sec:sparse}

\begin{theorem}[sparse class separation]\label{thm:sparse}
Let $K$ have at most one vertex of positive even degree, and let $\cA$ be a family of packed sets of $K$,
possibly overlapping, with union $X$, such that every vertex of $K$ has at most one neighbour in $X$. Then
$K$ has a Lov\'asz decomposition meeting every $A\in\cA$ at most once per path.
\end{theorem}

If $X$ is packed, every vertex has at most one neighbour in $X$, so Theorem~\ref{thm:sparse} with
$\cA=\{X\}$ is Theorem~\ref{thm:sep}.

\begin{proof}
Induction on $\nu(K)=2|E(K)|+[\text{no vertex of }K\text{ has positive even degree}]$, following
Lov\'asz's cases. The hypotheses pass to subgraphs, and to the graph $H(K;x,u)$ when $x\notin X$: the new
leaf is not in $X$, and $u$ trades the neighbour $x\notin X$ for $\ell\notin X$. Distances do not decrease,
so the sets of $\cA$ stay packed. The key observation, as in~\cite{RD}, is that a Fan transfer with
target outside $X$ preserves the conclusion, since by Lemma~\ref{lem:fan} only one path changes its vertex
set, by the target.

\emph{Case I: a vertex $z$ of positive even degree.} It has at least two neighbours and at most one in
$X$, so some $u\in N(z)\setminus X$. In $K-uz$ the only possible positive even vertex is $u$. By induction it
has a decomposition $\cD'$ as required. Every neighbour of $u$ in $K-uz$ is odd and has one end, so
$Z(u)=\emptyset<1=b_{\cD'}(z)$, and Lemma~\ref{lem:fan} restores $uz$ with target $u\notin X$.

\emph{Case II: no positive even vertex, and some $x\notin X$ with $d(x)\ge3$.} Take $u\in N(x)$ and
$H=H(K;x,u)$, in which $x$ is the only positive even vertex. By induction $H$ has a decomposition $\cD_H$
as required. Delete the leaf: the resulting decomposition $\cD'$ of $K-xu$ has $b_{\cD'}(u)\in\{0,2\}$ and
the right ends elsewhere, and its paths lie in those of $\cD_H$. If $b_{\cD'}(u)=0$ add the path $ux$, which
meets $X$ at most in $u$ because $x\notin X$. If $b_{\cD'}(u)=2$, every neighbour of $x$ in $K-xu$ has an
end, and Lemma~\ref{lem:fan} with $t=u$ and target $x\notin X$ restores $xu$.

\emph{Case III: no positive even vertex, every vertex outside $X$ of degree at most one, and some vertex
$v\notin X$ of degree one.} Let $vw$ be its edge. $K-vw$ has at most the positive even vertex $w$; decompose
it by induction and add the path $vw$.

\emph{Case IV: every vertex outside $X$ is isolated.} Then every edge has both ends in $X$, and since every
vertex has at most one neighbour in $X$, every degree is at most one. Every path of any Lov\'asz
decomposition is a single edge $ab$, which meets a packed set at most once because $a\sim b$.
\end{proof}

\begin{corollary}\label{cor:sparse}
Let $G$ and $N$ be as in Proposition~\ref{prop:reduction}, with no edge of $N$ in a triangle, and suppose
every vertex has at most one neighbour in $\Ev$ in the graph $F=G-N$. Then every orientation of $N$ is
realisable. Moreover any family of packed subsets of $\Ev$ (for example several pairings at distance three)
can be separated simultaneously.
\end{corollary}

For Gallai's bound alone Corollary~\ref{cor:sparse} is not new: the even subgraph has maximum degree at most
two and Fan's theorem applies. It is new as a statement about all orientations and about overlapping
families. In Theorem~\ref{thm:sparse} the Fan targets are always outside $X$; when every vertex lies in a
class this is impossible, which is the obstruction of Section~\ref{sec:caseI}.

\section{Evidence and limits (computational)}\label{sec:limits}

Nothing in Sections~\ref{sec:intro}--\ref{sec:sparse} depends on this section. The computations used the
exact oracle \texttt{tools/route\_ext.c} (query \texttt{S}), a dedicated checker for cubic graphs, and a
randomised constructive search; scripts are in \texttt{manuscripts/multi-sep/scripts}, and the log is
experiment E21 of \texttt{research/experiments.md}.

\subsection*{MSEP}
The cubic checker chooses an end edge at every vertex by depth-first search, keeping path components
with union--find and vertex masks, and was cross-checked against an independent enumerator on all cubic
graphs of order ten (285 queries, including 12 failures at distance two).
\begin{center}
\begin{tabular}{rrrl}
\toprule
order & connected cubic graphs & perfect pairings at distance $\ge3$ & separating decomposition\\
\midrule
10 & 19 & 54 (all) & all\\
12 & 85 & 2{,}846 (all) & all\\
14 & 509 & 269{,}966 (all) & all\\
16 & 4{,}060 & 36{,}024{,}731 (all) & all\\
18 & 41{,}301 & 8{,}260{,}200 (200 random per graph) & all\\
\bottomrule
\end{tabular}
\end{center}
For random graphs with all degrees odd (orders $12$--$20$, edge densities up to $0.5$, maximal pairings at
distance $\ge3$) no failure occurred. Instances left undecided by the budgeted oracle were all certified
by the randomised search; so were random $5$-regular graphs of orders $30$ and $40$ with perfect pairings.
In the hardest instances we found, random runs of Lov\'asz's induction separate in under $7\%$ of cases
(cubic) or almost never (dense random graphs), and a separating decomposition still exists.

\subsection*{CSEP}
No failure: classes of size three (399 instances, cubic, orders $12$--$16$); greedy maximal classes (orders
$14$, $18$); and \emph{minimum} proper colourings of the square (cubic orders $10$, $12$, $14$: $57$,
$170$, $150$ instances, with $4$ to $10$ colours; random graphs with all degrees odd, orders $10$, $12$:
$100$ each). Unions of two perfect pairings (each vertex in two pairs) never failed (626 instances); unions
of three can fail (2 of 228 at order $12$), so Remark~\ref{rem:overlap} is sharp in this sense.

\subsection*{The Case~I lemma}
Conjecture~\ref{conj:EE} held in every decided instance: $869$ pairs $(z,u)$ without classes and $788$
with pairings (random graphs, orders $8$--$15$; $51$ undecided within the search budget), and $2{,}272$ Case~I configurations obtained from cubic graphs with
perfect pairings by one Case~II step (orders $12$--$22$), each for every neighbour $u$.

\subsection*{What does not work}
\begin{itemize}
\item \emph{The plain hypothesis.} In about half of the Case~I configurations tested, no neighbour $u$ of
$z$ is such that every separating decomposition of $K-uz$ yields a separating decomposition of $K$ by
Fan's lemma. Most failing chains are trivial: the end path of $z$ contains the partner of $u$. So the
inductive hypothesis must be strengthened.
\item \emph{Strengthenings that are false.} (a) End paths avoid the partners of all neighbours: fails in
182 of 228 and 209 of 398 instances (cubic, orders 12, 14). (b) One forbidden class for every vertex: fails
in 114 of 199. (c) The even vertex as the centre of a path of length two: fails in 19 of 500. (d) Every
neighbour of the even vertex ends a path through it: fails in 39 of 302. (e) The end path inside the
closed neighbourhood: fails, already when the vertex is a leaf at the even vertex. (f) For cubic graphs,
decompositions into paths of length three built from a perfect matching and an oriented $2$-factor
(Kotzig's construction~\cite{Kot57}) do not always suffice: 2 failures in 228 at order 12.
\item \emph{Certified moves.} An induction using only Theorem~\ref{thm:caseII}, Lemma~\ref{lem:caseI} and
Lemma~\ref{lem:rotation} (exhaustive search over all choices) fails on small graphs: Case~I/II
certificates fail for 13 of 24 instances at order 10, for 3 of 63 triangle-free and 3 of 15 girth-five
instances at order 12, but for none of 108 girth-five instances at order 14. Rotation fails for 1 of 24 at
order 10 and 14 of 152 at order 12. On random $5$-regular graphs of order $30$ a randomised search for
Case~I/II certificates found none in ten instances, although separating decompositions exist. Local certificates need room at distance three, which dense or small graphs
do not have; a proof of MSEP must use information that is not local.
\item \emph{Several prescribed end edges.} Prescribing the terminal edge at every vertex of a packed set of
up to four vertices never failed ($540$ instances); compare Question~\ref{q:pe}.
\end{itemize}

\section{Formalisation}\label{sec:lean}

The following are proved in Lean~4 with Mathlib in \texttt{lean/GallaiLean/MultiSep.lean}, sorry-free; each
depends only on \texttt{propext}, \texttt{Classical.choice} and \texttt{Quot.sound}. They build on
\texttt{Separation.lean} (\texttt{separation\_set}, Theorem~\ref{thm:sep}) and \texttt{FanSupport.lean}.
\begin{itemize}
\item \texttt{PathDecomp.pushOn}: a decomposition pushed along a vertex map injective on every path and on
edges. \texttt{leaf\_merge}: Lemma~\ref{lem:merge}.
\item \texttt{lovasz\_endEdge\_set}, \texttt{lovasz\_endEdge}: Theorem~\ref{thm:endedge}.
\item \texttt{ClassFam}, \texttt{ClassSep}, \texttt{caseII\_lift}, \texttt{classSep\_allOdd}:
Theorem~\ref{thm:caseII}.
\item \texttt{classSep\_caseI}: Lemma~\ref{lem:caseI}.
\item \texttt{SparseFam}, \texttt{separation\_sparse}: Theorem~\ref{thm:sparse}, by the same case structure as
\texttt{separation\_set}, with a new Case~IV.
\item \texttt{realisable\_of\_separating}, \texttt{realisable\_of\_lovasz\_separating}:
Proposition~\ref{prop:reduction}(c). If $\sigma$ pairs the even vertices of $G$ by edges of $G$
(\texttt{IsEvenPairing}), $F$ is $G$ minus the pairing edges, and $F$ has a Lov\'asz decomposition in
which no path contains both $v$ and $\sigma(v)$, then every transversal of the pairing is realisable
(\texttt{Realisable}). The proof appends single-edge routes with \texttt{append\_routes} of
\texttt{RouteSystem.lean} (route-descent).
\end{itemize}
Not formalised: Propositions~\ref{prop:reduction}(a),(b), \ref{prop:cover} and~\ref{prop:EE}, and
Lemma~\ref{lem:rotation}.

\section{Open problems}\label{sec:open}

\begin{question}\label{q:EE}
Prove Conjecture~\ref{conj:EE}, even without classes: if $z$ is the only vertex of positive even degree
and $u\in N(z)$, is $uz$ the terminal edge at $u$ of some Lov\'asz decomposition? By
Proposition~\ref{prop:EE} this is $L_1$ at the even vertex, for the avoided set $\{u\}$.
\end{question}

\begin{question}\label{q:cubic}
Prove Conjecture~\ref{conj:msep} for cubic graphs (all orientations for $4$-regular graphs whose
perfect matching avoids triangles), or for two pairs.
\end{question}

\begin{question}[prescribed end edges on a packed set]\label{q:pe}
If all degrees are odd, $X$ is packed and an edge $xu_x$ is chosen at every $x\in X$, is there a Lov\'asz
decomposition in which every $xu_x$ is the terminal edge at $x$? Theorem~\ref{thm:endedge} is the case
$|X|=1$.
\end{question}

\begin{question}\label{q:csep}
Is Conjecture~\ref{conj:csep} true? Equivalently: for every proper colouring of the square of a graph
with at most one vertex of positive even degree, is there a Lov\'asz decomposition into rainbow paths?
\end{question}

\begin{thebibliography}{99}
\bibitem{BP19} M.~Bonamy and T.~J.~Perrett, Gallai's path decomposition conjecture for graphs of small
maximum degree, Discrete Math.\ 342 (2019) 1293--1299.
\bibitem{BS21} F.~Botler and M.~Sambinelli, Towards Gallai's path decomposition conjecture, J.~Graph
Theory 97 (2021) 161--184.
\bibitem{CW25} Y.~Chu and J.~Wang, Path decompositions of Eulerian graphs, arXiv:2510.12806 (2025).
\bibitem{Fan05} G.~Fan, Path decompositions and Gallai's conjecture, J.~Combin.\ Theory Ser.~B 93 (2005)
117--125.
\bibitem{FK} O.~Favaron and M.~Kouider, Path partitions and cycle partitions of Eulerian graphs of maximum
degree~4, Studia Sci.\ Math.\ Hungar.\ 23 (1988) 237--244.
\bibitem{Kot57} A.~Kotzig, Z te\'orie kone\v{c}n\'ych pravideln\'ych grafov tretieho a \v{s}tvrt\'eho
stup\v{n}a, \v{C}asopis Pest.\ Mat.\ 82 (1957) 76--92.
\bibitem{L68} L.~Lov\'asz, On covering of graphs, in: Theory of Graphs (Proc.\ Colloq.\ Tihany, 1966),
Academic Press, New York, 1968, 231--236.
\bibitem{Pyb96} L.~Pyber, Covering the edges of a connected graph by paths, J.~Combin.\ Theory Ser.~B 66
(1996) 152--159.
\bibitem{RD} J.~Williams, Route descent: transversal realisability, minimum $T$-joins, and Lov\'asz
decompositions that separate two vertices, manuscript (\texttt{manuscripts/route-descent}), 2026.
\end{thebibliography}
\end{document}
